\section{问题一模型建立与求解}

\subsection{问题分析}
% 明确问题一的输入、输出、评价指标和建模难点。
在此给出问题一的技术路线。图~\ref{fig:problem1-flow} 可替换为实际算法流程图。
\figureplaceholder{问题一建模流程图}{fig:problem1-flow}

\subsection{数据处理}
% 说明本问题实际使用的样本、特征、训练集/测试集划分及防止数据泄漏的措施。
在此引用第五节的数据预处理结果，并说明问题一的专用处理步骤。

\subsection{模型建立}
% 给出变量、目标函数、约束和参数含义，而不是只列算法名称。
用占位表达式表示预测或评价模型：
\begin{equation}
  \hat{y}_i=f(\boldsymbol{x}_i;\boldsymbol{\theta}),
  \qquad
  \boldsymbol{\theta}^{*}=\arg\min_{\boldsymbol{\theta}}L(\boldsymbol{\theta}).
  \label{eq:model-placeholder}
\end{equation}

\subsection{模型求解}
% 说明算法步骤、参数设置、软件环境、随机种子和停止条件。
在此描述训练、搜索或优化过程，并把完整可运行程序放入附录与支撑材料。

\subsection{结果分析}
% 报告点估计、误差、置信区间或分类指标，并解释结果的实际含义。
\begin{table}[H]
  \centering
  \caption{问题一结果占位表}
  \label{tab:problem1-result}
  \begin{tabularx}{0.78\textwidth}{cXXX}
    \toprule
    模型 & 训练集指标 & 验证集指标 & 备注 \\
    \midrule
    基准模型 & 【填写】 & 【填写】 & 用于比较 \\
    改进模型 & 【填写】 & 【填写】 & 最终模型 \\
    \bottomrule
  \end{tabularx}
\end{table}
